\documentclass[10pt,letterpaper]{article}

% DeepRare-inspired single-column manuscript template.
% Compile with: latexmk -xelatex main.tex

\usepackage[letterpaper,top=0.72in,bottom=0.70in,left=0.92in,right=0.92in]{geometry}
\usepackage{microtype}
\usepackage{amsmath,amssymb}
\usepackage{graphicx}
\usepackage{booktabs}
\usepackage{enumitem}
\usepackage{xcolor}
\usepackage{caption}
\usepackage{fancyhdr}
\usepackage{titlesec}
\usepackage[backend=biber,style=numeric,sorting=none,maxbibnames=99]{biblatex}
\usepackage[colorlinks=true,citecolor=blue!70!black,linkcolor=blue!55!black,urlcolor=blue!70!black]{hyperref}

\addbibresource{references.bib}

\setlength{\parindent}{0pt}
\setlength{\parskip}{0.42em}
\setlength{\emergencystretch}{2em}
\setlist[itemize]{leftmargin=2.2em,itemsep=0.25em,topsep=0.3em}
\captionsetup{font=small,labelfont=bf}
\titleformat{\section}{\large\bfseries}{\thesection}{0.65em}{}
\titleformat{\subsection}{\normalsize\bfseries}{\thesubsection}{0.65em}{}
\titlespacing*{\section}{0pt}{1.0em}{0.25em}
\titlespacing*{\subsection}{0pt}{0.75em}{0.15em}

\pagestyle{fancy}
\fancyhf{}
\renewcommand{\headrulewidth}{0pt}
\fancyfoot[R]{\small $\vert$\,\thepage}

\newcommand{\systemname}{\textbf{Medical Skill Hub}}

\title{\vspace{-1.8em}\bfseries
An Agentic System for Evidence-Grounded\\
Digestive Disease Diagnosis}

\author{
First Author$^{1,*}$, Second Author$^{1,*}$, Third Author$^{2}$,\\
Senior Author$^{1,\dagger}$ and Corresponding Author$^{2,\dagger}$\\[0.55em]
\small $^{1}$Department or School, University, City, Country\\
\small $^{2}$Hospital or Research Institute, City, Country\\[0.35em]
\small $^{*}$Equal contributions \qquad $^{\dagger}$Corresponding authors\\
\small Correspondence: \href{mailto:author@example.com}{author@example.com}
}
\date{}

\begin{document}

\maketitle
\thispagestyle{fancy}
\vspace{-1.5em}

% DeepRare uses an unlabelled summary rather than a conventional Abstract heading.
Accurate differential diagnosis in gastroenterology requires the integration of
longitudinal clinical narratives, similar cases, current biomedical literature,
and specialty guidelines. Here we present \systemname, an agentic clinical
decision-support system that produces ranked diagnostic candidates with
ICD-10-CM codes and traceable supporting evidence. The system combines an
initial language-model assessment with hybrid similar-case retrieval, PubMed
search, and locally compiled guideline skills. We evaluate the system across
diverse digestive diseases using stage-specific and final Top-$k$ recall. This
template paragraph should be replaced with a concise summary of the clinical
motivation, methods, principal numerical findings, and implications of the work.

\section{Main}

Digestive diseases present substantial diagnostic challenges because symptoms,
laboratory abnormalities, imaging findings, and endoscopic observations often
overlap across conditions. Clinical decision-support systems must therefore
combine broad medical knowledge with patient-specific evidence while preserving
traceability and clearly separating observed facts from retrieved information.

Recent progress in large language models and retrieval-augmented generation has
enabled systems that synthesize heterogeneous evidence sources
\cite{singhal2023clinical,lewis2020retrieval}. However, diagnostic performance
alone is insufficient for medical use: each candidate should be supported by
verifiable evidence, and external knowledge must not be presented as an observed
patient fact.

Here, we introduce \systemname, a multi-stage framework for digestive disease
diagnosis. Given a complete case narrative, the system generates initial
hypotheses, retrieves similar cases, plans literature searches, consults local
guideline skills, and produces five unique ranked diagnoses with ICD-10-CM
codes. A diagnostic judgement stage determines whether a second evidence-
gathering round is needed.

Our main contributions are:

\begin{itemize}
  \item an evidence-grounded, multi-agent workflow covering the full spectrum of
        digestive diseases rather than a single condition;
  \item hybrid case retrieval that combines lexical and dense representations;
  \item traceable use of PubMed literature and locally compiled clinical
        guidelines; and
  \item stage-wise evaluation of diagnostic recall and evidence usage.
\end{itemize}

\section{Results}

\subsection{System Overview}

Figure~\ref{fig:overview} summarizes the diagnostic workflow. The preprocessing
agent extracts positive clinical features and proposes initial hypotheses. A
retrieval module ranks local cases using BM25 and dense similarity, fuses the
rankings, and applies a medical cross-encoder. Search-planning, literature, and
guideline agents then collect evidence for the final diagnostic agent.

\begin{figure}[htbp]
  \centering
  % Replace this framed box with:
  % \includegraphics[width=\textwidth]{figures/system-overview.pdf}
  \fbox{\rule{0pt}{2.15in}\rule{0.94\textwidth}{0pt}}
  \caption{\textbf{Overview of the Medical Skill Hub workflow.}
  Replace the placeholder with a vector PDF or high-resolution image. Describe
  each panel in order and define all abbreviations.}
  \label{fig:overview}
\end{figure}

\subsection{Evaluation Settings}

Describe the study population, inclusion criteria, case sources, comparison
methods, model settings, and primary endpoints. Keep training, validation, and
test cohorts clearly separated. For clinical datasets such as MIMIC-IV
\cite{johnson2023mimic}, report governance and data-access requirements.

The primary metrics are category- and subcategory-level Recall@1, Recall@3, and
Recall@5. Confidence intervals and statistical tests should be reported where
appropriate. Table~\ref{tab:main-results} gives a compact example layout.

\begin{table}[htbp]
  \centering
  \caption{\textbf{Main diagnostic performance.} Replace all values with study
  results and report uncertainty where applicable.}
  \label{tab:main-results}
  \begin{tabular}{lccc}
    \toprule
    Method & Recall@1 & Recall@3 & Recall@5 \\
    \midrule
    Initial LLM hypotheses & -- & -- & -- \\
    Similar-case retrieval & -- & -- & -- \\
    Medical Skill Hub      & -- & -- & -- \\
    \bottomrule
  \end{tabular}
\end{table}

\subsection{Overall Diagnostic Performance}

Report the principal comparison first, including cohort size, effect magnitude,
and uncertainty. Follow with results across disease categories, care settings,
or case-complexity strata. Avoid restating every value shown in a table.

\subsection{Contribution of Evidence Sources}

Quantify the contribution of similar cases, biomedical literature, and local
guideline skills. Distinguish evidence availability from evidence correctness,
and document cases in which an external source could not be retrieved.

\subsection{Ablation and Error Analysis}

Evaluate key pipeline components and summarize clinically meaningful failure
modes. Suggested groups include insufficient case information, competing
diagnoses with overlapping presentations, coding granularity errors, retrieval
failures, and unsupported model inferences.

\section{Discussion}

Interpret the findings in relation to the clinical problem and prior work.
Discuss why particular components improved or failed to improve performance,
how traceable evidence may support clinician review, and which conclusions are
supported by the reported experiments. Avoid implying autonomous clinical use.

\section{Limitations}

Discuss dataset representativeness, retrospective design, coding quality,
possible information leakage, dependence on external models and databases,
guideline coverage, and the need for prospective evaluation. State explicitly
that the system supports research and technical demonstration and does not
replace clinician diagnosis, treatment decisions, or in-person assessment.

\section{Acknowledgments}

List funding sources, institutional support, data providers, and individuals who
contributed but do not meet authorship criteria. Include grant numbers where
required.

\section{Author Contributions}

Use the CRediT taxonomy where possible. Example: F.A. and S.A. conceived the
study; F.A. implemented the system; F.A. and T.A. performed experiments; all
authors analyzed results, revised the manuscript, and approved the final version.

\section{Competing Interests}

The authors declare no competing interests. Replace this statement if any
financial or non-financial interests must be disclosed.

\section{Methods}

\subsection{Problem Formulation}

We formulated digestive-disease diagnosis as principal-diagnosis ranking from a
complete free-text clinical record. For a case $x$, the system returns an ordered
set $\widehat{Y}_{5}=\{(d_j,c_j,q_j,e_j,a_j)\}_{j=1}^{5}$, where $d_j$ is a
unique ICD-10-CM code, $c_j$ is its canonical English name, $q_j$ is an integer
confidence score from 0 to 100, $e_j$ contains case-grounded supporting evidence,
and $a_j$ contains recommended next examinations or management directions. The
first-ranked item is intended to represent the condition chiefly responsible for
the admission or the main condition evaluated and treated during the encounter.
Codes are normalized to uppercase and stored without decimal points. Three-
character categories are used only when a more specific supported subcategory is
not available.

We enforced a strict provenance boundary throughout the pipeline. The original
case text was the only permissible source of facts about the current patient.
Initial hypotheses, retrieved cases, PubMed abstracts, guideline statements, and
outputs from previous rounds were treated as candidate-generation signals or
external knowledge, not as observations in the current patient. Consequently,
the system could use external evidence to interpret a documented finding but
could not import a symptom, test result, or diagnosis from another source into
the patient's record.

\subsection{Overall Diagnostic Workflow}

The pipeline comprised six linked stages. First, two independent preprocessing
agents generated an initial differential diagnosis and extracted structured
positive findings from the original record. Second, the structured findings were
used to retrieve diagnostically similar cases from a local case bank. Third, the
five initial hypotheses and the leading similar-case diagnoses were merged and
ranked to form a candidate set, and a search-planning agent generated PubMed
queries for that set. Fourth, PubMed retrieval and local guideline-skill search
were executed concurrently. Fifth, a diagnosis agent ranked the supplied
candidates using the patient record and the retrieved external evidence. Finally,
an evidence-sufficiency agent determined whether a second, targeted literature
retrieval round was warranted.

The initial preprocessing and similar-case retrieval were performed once. The
system allowed at most two diagnostic rounds. When the first-round evidence was
sufficient, the pipeline returned immediately. Otherwise, the second round
retained the original candidate set, similar-case results, and guideline results,
but regenerated focused PubMed queries and accumulated newly selected abstract
sections before revising the final ranking. This bounded design prevented an
open-ended agent loop while preserving one opportunity to resolve a specific
evidence gap.

\subsection{Independent Clinical Preprocessing}

Two agents processed the same original case in parallel and were not shown each
other's outputs. The hypothesis agent independently produced exactly five unique
principal-diagnosis candidates, each represented by a complete ICD-10-CM code and
canonical English name. Chronic comorbidities, incidental findings, symptoms,
and secondary complications were not used merely to fill lower ranks.

The feature-extraction agent organized explicitly documented positive findings
into five fields: present illness history, past medical history, physical
examination, family history, and pertinent laboratory, imaging, endoscopic,
pathological, and microbiological results. Each finding was represented as a
short English phrase while retaining clinically relevant duration, severity,
anatomical location, measurements, and trends. Denied symptoms, normal or
negative examinations, planned tests, and inferred findings were excluded. Empty
fields were preserved as empty lists. These field-specific representations were
used only as queries for similar-case and guideline retrieval; all subsequent
diagnostic claims remained anchored to the full original record.

\subsection{Similar-Case Retrieval}

The local case bank stored one admission identifier, discharge diagnosis,
ICD-10-CM code, and the same five clinical fields for each case. Non-empty fields
were segmented into non-overlapping chunks of at most 510 tokenizer tokens. For
the current patient, all items within each non-empty feature field were
concatenated to create one field-specific query. Retrieval was conducted
independently within the corresponding field of the case bank, avoiding, for
example, direct comparison of a laboratory query with a family-history section.

We used parallel lexical and dense retrieval branches. The lexical branch applied
BM25 \cite{robertson1995okapi} after lower-casing and tokenizing text with the
scientific spaCy tokenizer, and retained up to 50 positive-scoring sections per
field. The dense branch encoded each query and case section with BGE-M3
\cite{chen2024bgem3}. We used the normalized embedding of the first output token,
truncated model inputs to 1,024 tokens, and ranked sections by the dot product of
L2-normalized vectors, which is equivalent to cosine similarity. Up to 50 dense
candidates were retained per field. Corpus embeddings and BM25 indices were
cached and invalidated when the source case-bank fingerprint changed.

Section scores were aggregated at the admission level. For each field and
retrieval branch, the case score was the best section score plus 0.2 times the
second-best score when a second matching section existed. The five feature fields
were assigned weights of 1.0, 0.5, 0.5, 0.2, and 1.0 for present illness, past
medical history, physical examination, family history, and pertinent results,
respectively. We combined ranks using reciprocal rank fusion (RRF)
\cite{cormack2009rrf}. Let $\mathcal{F}$ denote the non-empty patient fields,
$\bar{w}_f=w_f/\sum_{g\in\mathcal{F}}w_g$ the normalized field weight, and
$r_{m,f}(i)$ the rank of case $i$ for modality $m$ and field $f$. The modality-
specific fusion score was

\begin{equation}
  S_m(i)=\sum_{f\in\mathcal{F}}
  \frac{\bar{w}_f}{60+r_{m,f}(i)},
\end{equation}

where a term was omitted when the case did not occur in that ranked list. The
hybrid score combined the lexical and dense branches as

\begin{equation}
  S_{\mathrm{hybrid}}(i)=\sum_{f\in\mathcal{F}}\bar{w}_f
  \left[
  \frac{0.4}{60+r_{\mathrm{BM25},f}(i)}+
  \frac{0.6}{60+r_{\mathrm{dense},f}(i)}
  \right].
\end{equation}

The fused ranking was deduplicated by ICD-10-CM code and truncated to 20 disease
candidates. When more than five candidates were available, an LLM reranker was
given anonymous candidate identifiers, diagnosis codes and names, BM25, dense,
and RRF ranks, and the two highest-scoring fused sections. It selected exactly
five unique candidates according to their fit with the current patient's
structured findings. Similar-case text remained explicitly marked as external
case evidence. If reranking failed validation or timed out after 120 seconds, the
top five RRF candidates were used; when five or fewer fused candidates were
available, their RRF order was retained directly. The pipeline also preserved
the distinct top-five BM25, dense, RRF, and LLM-reranked outputs for analysis.

\subsection{Candidate Integration and Search Planning}

Candidate construction began with the five independently generated hypotheses,
followed by up to five LLM-reranked similar-case diagnoses, with the top RRF cases
used as fallback. Exact normalized ICD-10-CM duplicates were removed while
preserving the first occurrence, yielding at most ten candidates. When at least
two candidates were present, a separate planning reranker ordered all supplied
codes against the original case. Exact, four-character, and three-character
agreement between the initial-LLM and similar-case sources served as secondary
signals and could not override contradictory patient findings. The reranker was
required to return an exact permutation of the input codes; otherwise, the
initial merge order was retained.

The search-planning agent copied this ranked candidate set without adding,
removing, or renaming a diagnosis. In the initial round, it generated five to ten
concise PubMed-oriented keyword queries whose collective coverage included every
candidate and the most discriminative documented patient features. Queries could
target diagnostic criteria, anatomical distributions, imaging, endoscopy,
histopathology, procedural complications, or pairwise differential diagnosis.
If a second round was requested, the agent instead generated one to five focused
queries restricted to the diagnoses and knowledge gaps specified by the
evidence-sufficiency judgement.

\subsection{PubMed Retrieval and Evidence Selection}

Each planned query was submitted to the PubMed E-utilities ESearch endpoint, and
up to three identifiers were retrieved. The corresponding records were then
downloaded from EFetch in XML format. Requests included the configured tool name,
email address, and API key when available. A process-wide rate limiter enforced
the configured request frequency; HTTP 429 and server errors were retried with
exponential delay and jitter. Queries from the same planning result were
executed concurrently.

The parser retained only records with a numeric PMID, non-empty title, and at
least one abstract section. Structured abstracts were preserved as separate,
one-indexed sections rather than collapsed into a single text block. A knowledge-
search agent received all retrieved titles and sections and returned only exact
$(\mathrm{PMID},\mathrm{section\ index})$ pairs judged relevant to one or more
planned queries. Repeated retrieval of a PMID across queries was treated as a
relevance signal but was not sufficient for selection. The agent could neither
rewrite an abstract nor invent a PMID or section index.

\subsection{Local Guideline-Skill Retrieval}

Clinical guidelines were prepared offline as local Agent Skills. Each skill
contained disease-scope metadata, workflow instructions, a recommendation index,
the converted guideline full text, and a local search utility. At run time, the
guideline orchestrator compared the candidate diagnoses with the skill catalog.
A skill was selected directly when its primary disease matched a candidate or
when its broader scope clearly covered that candidate. A narrower skill was
eligible only if all of its required qualifiers, such as subtype, complication,
stage, hereditary status, or procedure context, were present in the candidate.
Shared organ systems, symptoms, or treatments alone were not sufficient.

For every directly matched skill, the system performed one non-recursive forward
expansion through differential diseases explicitly declared in that skill's
metadata. An additional skill was selected only when its primary scope directly
matched one of these declared differentials; reverse, multi-hop, and generic
category expansion were prohibited. Selected skills were executed concurrently,
with at most ten active executions. Each executor ran in a read-only sandbox,
loaded exactly one selected skill, read its complete instructions, used the
recommendation index to locate relevant passages, and verified those passages
against the guideline full text.

Each executor returned the disease assessed, verified evidence, a concise
patient-to-guideline comparison, and one of four applicability labels:
\emph{supported}, \emph{possible}, \emph{insufficient evidence}, or
\emph{not supported}. Empty guideline evidence necessarily produced the
insufficient-evidence label. Before final diagnosis, directly matched skills were
retained when labelled supported or possible, whereas differentially expanded
skills were retained only when labelled supported. This stricter rule prevented
a one-hop expansion from entering the final evidence set solely because it was
compatible with the patient.

\subsection{Candidate-Constrained Final Diagnosis}

The final diagnosis agent received the original case, the candidate diagnoses,
selected guideline assessments, and selected PubMed sections. Each candidate was
annotated with its planning rank, source membership, initial-LLM rank, similar-
case rank, and exact, four-character, and three-character cross-source agreement.
Guideline evidence was numbered first, followed by PubMed evidence. Candidate
metadata and numbered material were treated as ranking priors and external
knowledge, respectively, rather than patient facts.

The agent was constrained to return exactly five unique codes from the supplied
candidate set and to preserve each corresponding canonical name. Rank 1 was
optimized for principal-diagnosis precision, ranks 1--3 for the strongest
alternatives, and ranks 1--5 for plausible diagnostic coverage. The first initial
LLM hypothesis was retained at rank 1 unless explicit patient findings favored
another supplied candidate and cross-source evidence also supported that change.
Lower ranks could add three- or four-character ICD diversity only when clinical
support was otherwise comparable; diversity alone could not displace a stronger
candidate. Unsupported etiological, anatomical, complication, or severity
specificity was prohibited.

For each diagnosis, the output included an independent integer confidence from 0
to 100, supporting findings copied or concisely derived from the current record,
and recommended next steps. Every planning candidate omitted from the top five
had to be returned once with a patient-grounded reason for exclusion. Programmatic
validation rejected new codes, altered names, duplicate diagnoses, ranks other
than 1--5, missing excluded candidates, and empty exclusion reasons. One
correction attempt was allowed after a validation error. Finally, only external
evidence actually cited in a diagnosis, exclusion reason, or next step was kept,
and citation numbers were compactly reassigned in order of first appearance.

\subsection{Evidence-Sufficiency Judgement and Iterative Refinement}

After each final diagnosis, a separate judgement agent reviewed the patient
record, search plan, retrieved PubMed and guideline results, and complete ranked
output. It requested a second round only when a specific unresolved distinction
could materially change the principal-diagnosis ranking or ICD-10-CM selection,
the available external evidence did not address that distinction, and another
literature search could realistically resolve it. Uncertainty requiring a new
patient-specific test did not qualify as an external evidence gap.

When further retrieval was requested, the judgement output had to identify
unique focus codes from the current top five, specific evidence gaps, and PubMed
query directions. The search-planning agent used this feedback to revise the
queries but retained the original candidate set. During round 2, only PubMed
search was rerun; similar-case retrieval and guideline-skill execution were not
repeated. Newly selected PubMed sections were merged with round-1 evidence and
deduplicated by PMID and section index. The diagnosis agent then reassessed all
candidates using the accumulated evidence and the previous judgement.

If the second judgement still requested more evidence, no third retrieval round
was initiated. Instead, the system performed one corrective final-diagnosis pass
using the accumulated evidence, explicitly preserving unresolved uncertainty,
avoiding unsupported code specificity, and reducing confidence when appropriate.
The last completed diagnosis was returned as the pipeline output.

\subsection{Failure Handling and Implementation}

The orchestration was implemented in Python 3.10 with the OpenAI Agents SDK, and
all inter-agent outputs were validated with Pydantic schemas. The diagnosis model
was configurable between the OpenAI Responses API, a DeepSeek OpenAI-compatible
endpoint, and locally served Qwen models. Native structured output was used when
supported; compatible chat-completion models were instructed to return JSON
objects. Model temperature was set to zero where supported.

The external PubMed and guideline branches were run with independent exception
handling. Failure of either branch produced an empty stage result and a failure
reason while allowing diagnosis to continue with the remaining evidence.
Similar-case retrieval failure was handled in the same way. If search planning
failed, the merged candidate list was preserved and PubMed queries were left
empty. These fallbacks affected evidence availability but did not relax the final
five-code schema or the constraint that patient facts originate only from the
case record.

\subsection{Evaluation}

Batch evaluation used the recorded principal ICD-10-CM code as the reference
standard. We evaluated seven outputs: the initial LLM hypotheses; BM25, dense,
RRF, and LLM-reranked similar cases; the merged search-planning candidates; and
the final diagnoses. Before comparison, reference and predicted codes were
upper-cased and stripped of decimal points. Disease-category agreement required
the first three characters to match, whereas subcategory agreement required the
first four characters to match.

For a dataset of $N$ cases and an ICD prefix length $p\in\{3,4\}$, Recall@$k$
was computed as

\begin{equation}
  \operatorname{Recall@}k^{(p)} = \frac{1}{N}\sum_{i=1}^{N}
  \mathbb{I}\!\left(y_i^{(p)} \in \widehat{Y}_{i,k}^{(p)}\right),
\end{equation}

where $y_i^{(p)}$ is the length-$p$ prefix of the reference code and
$\widehat{Y}_{i,k}^{(p)}$ is the set of corresponding prefixes among the top $k$
predictions. Recall@1, Recall@3, and Recall@5 were reported for all stages;
Recall@20 was additionally available for the RRF candidate list, and Recall@10
for the search-planning list. Metrics were summarized for each diagnostic round
and for the final returned round. Guideline-skill usage rate and per-skill usage
counts were reported separately.

\section{Data Availability}

Describe which data can be shared, where they can be accessed, and any clinical
privacy or licensing restrictions.

\section{Code Availability}

Provide the repository URL, release identifier, license, environment setup, and
instructions needed to reproduce the reported experiments.

\printbibliography[title={References}]

\section{Figure Legends}

\textbf{Figure 1. Overview of the Medical Skill Hub workflow.}
Provide a self-contained legend, explain each panel, and define every acronym,
color, symbol, and statistical annotation.

\appendix
\section{Supplementary Information}

Place extended methods, prompts, additional analyses, dataset descriptions, and
case studies here, or move them to a separate supplementary manuscript if the
target journal requires separate submission files.

\end{document}
